\documentclass{article}
\usepackage{graphicx} % Required for inserting images

\usepackage{amsmath}
\usepackage{amssymb}
\newcommand{\bigp}[1]{\left(#1\right)}
\newcommand{\bigb}[1]{\left[#1\right]}
\newcommand{\summation}[1][0]{\sum_{n=#1}^\infty}
\newcommand{\summationk}[1][0]{\sum_{k=#1}^\infty}
\newcommand{\summationm}[1][0]{\sum_{m=#1}^\infty}

\title{Infinite Series 3}
\author{Siddharth K}
\date{January 2026}

\begin{document}

\maketitle

\section*{Problem 2}
Rewriting $D_N(t)$ in terms of exponentials:
$$D_N(t) = \frac{1}{2} + \sum_{n = 1}^N \cos(nt) = \frac{1}{2}\bigp{1 + \sum_{n = 1}^N (e^{int} + e^{-nt})}$$
$$= \frac{1}{2}\sum_{n = -N}^N e^{int}.$$
This is a geometric series:
$$a = e^{-iNt}, \quad r = e^{it},$$
$$a\sum_{n = 0}^{2N}r^n = \frac{a(r^{2N + 1} - 1)}{r - 1} = \frac{e^{-iNt}(e^{(2N + 1)t - 1})}{e^{it} - 1} = \frac{e^{i(N + 1)t} - e^{iNt}}{e^{it} - 1}.$$
Multiplying the numerator and denominator by $e^{-it/2}$:
$$= \frac{e^{i(N + 1/2)t} - e^{i(N + 1/2)t}}{e^{it/2} - e^{-it/2}} = \frac{2\sin((N + 1/2)t)}{2\sin(t/2)}.$$
Finally, multiplying by $1/2$:
$$D_N(t) = \frac{\sin((N + 1/2)t)}{2\sin(t/2)}.$$

\section*{Problem 3}
$$\int_0^\pi D_N(t)dt = \int_0^{\pi} \bigp{\frac{1}{2} + \sum_{n = 1}^N \cos(nt)}dt = \frac{\pi}{2} + \sum_{n = 1}^N \int_0^\pi \cos(nt)dt$$
$$= \frac{\pi}{2} + \sum_{n = 1}^N -\frac{1}{n}\sin(\pi n) = \frac{\pi}{2}.$$

\section*{Problem 5}
Because $f$ is even, the Fourier series has only cosine terms:
$$a_0 = \frac{1}{2\pi}\int_{-\pi}^\pi f(x)dx, \quad a_n = \frac{1}{\pi}\int_{-\pi}^\pi f(x)\cos(nx)dx.$$
$$a_0 = \frac{1}{2\pi}\int_{-\pi}^\pi \cos(ax)dx = \frac{1}{2\pi}\bigb{\frac{\sin(ax)}{a}}_{-\pi}^\pi = \frac{1}{2\pi}\frac{2\sin(a\pi)}{a\pi}.$$
To compute $a_n$, start with product-to-sum:
$$\cos(ax)\cos(nx) = \frac{1}{2}(\cos((a - n)x) + \cos((a + n)x)).$$
So
$$a_n = \int_{-\pi}^\pi \cos(ax)\cos(nx)dx = \frac{1}{2\pi}\int_{-\pi}^\pi (\cos((a - n)x) + \cos((a + n)x))dx$$
$$= \frac{1}{2\pi}\bigp{\frac{2\sin((a - n)\pi)}{a - n} + \frac{2\sin((a + n)\pi)}{a + n}} = \frac{1}{\pi}\bigp{\frac{\sin((a - n)\pi)}{a - n} + \frac{\sin((a + n)\pi)}{a + n}}.$$
Now, simplifying the sines:
$$\sin((a \pm n)\pi) = \sin(a\pi \pm n\pi) = \sin(a\pi)\cos(a\pi) \pm \cos(a\pi)\sin(a\pi) = \sin(a\pi)(-1)^n,$$
$$a_n = \frac{1}{\pi}\bigp{\frac{\sin(a\pi)(-1)^n}{a - n} + \frac{\sin(a\pi)(-1)^n}{a + n}} = \frac{\sin(a\pi)(-1)^n}{\pi}\bigp{\frac{1}{a - n} + \frac{1}{a + n}}.$$
$$= \frac{2a\sin(a\pi)(-1)^n}{\pi(a^2 - n^2)}.$$
So the full Fourier series is
$$\cos(ax) = \frac{\sin(a\pi)}{a\pi} + \summation[1]\bigp{\frac{2a\sin(a\pi)(-1)^n}{\pi(a^2 - n^2)}}\cos(nx).$$
Since $f$ is defined on $[-\pi, \pi)$ and extended periodically, we have that
$$f(\pi) = f(-\pi) = \cos(-a\pi) = \cos(a\pi).$$
Therefore, the Fourier series at $\pi$ converges to $\cos(a\pi)$. At $x = \pi$, $\cos(n\pi) = (-1)^n$, so
$$\cos(a\pi) = \frac{\sin(a\pi)}{a\pi} + \summation[1]\bigp{\frac{2a\sin(a\pi)(-1)^n}{\pi(a^2 - n^2)}}(-1)^n$$
$$= \frac{\sin(a\pi)}{\pi}\bigp{\frac{1}{a} + \summation[1] \frac{2a}{a^2 - n^2}}.$$
Since $a \notin \mathbb{Z}$, we can divide both sides by $\sin(a\pi)$:
$$\pi\cot(a\pi) = \frac{1}{a} + \summation[1] \frac{2a}{a^2 - n^2},$$
which is the cotangent sum identity.

\section*{Problem 9}
When the period is $p$, the frequency becomes $\frac{2\pi}{p}$, so the exponent becomes $e^{i2\pi x/p}$. Given any starting point $x_0$, we integrate over one period:
$$c_n = \frac{1}{p}\int_{x_0}^{x_0 + p} f(x)e^{-i2\pi n x/p}\,dx.$$
Similarly,
$$a_n = \frac{2}{p}\int_{x_0}^{x_0 + p} f(x)\cos\bigp{\frac{2\pi nx}{p}}dx, \quad b_n = \frac{2}{p}\int_{x_0}^{x_0 + p} f(x)\sin\bigp{\frac{2\pi nx}{p}}dx.$$

\section*{Problem 10}
\subsection*{Part (b)}
Let $u = \lambda t, dt = du/\lambda$:
$$\widehat{g}(\xi) = \int e^{-2\pi\xi t}f(\lambda t)dt = \frac{1}{\lambda}\int e^{-2\pi i(\xi/\lambda)u}f(u) = \frac{1}{\lambda}\widehat{f}(\xi/\lambda).$$

\subsection*{Part (c)}
Take the Fourier transform:
$$\widehat{f * g}(\xi) = \int e^{-2\pi i\zeta t}\bigp{\int f(x)g(t - x)dx}dt$$
$$= \int f(x)\bigp{\int e^{-2\pi i\zeta t}g(t - x)dt}dx.$$
Substitute $u = t - x$:
$$= \int f(x)e^{-2\pi i\xi x}\bigp{\int e^{-2\pi\xi u}g(u)du}dx = \widehat{f * g}(\xi) = \widehat{f}(\xi)\widehat{g}(\xi).$$

\subsection*{Part (d)}
$$\widehat{f'}(\xi) = \int e^{-2\pi i\zeta t}f'(t)dt = \bigb{e^{-2\pi i\xi t}f(t)}_{-\infty}^\infty + \int (2\pi i\xi)e^{-2\pi i\xi t}f(t) dt.$$
The boundary term vanishes because $f$ decays faster than any power, so
$$\widehat{f'}(\xi) = \int (2\pi i\xi)e^{-2\pi i\xi t}f(t) dt = 2\pi i\xi \widehat{f}(\xi).$$

\section*{Problem 11}
Since $f$ is only nonzero on $(-\pi, \pi)$, the integral becomes
$$\widehat{f}(\xi) = \int_{-\pi}^\pi e^{-2\pi i\xi x}dx = \bigb{\frac{e^{-2\pi i\xi x}}{-2\pi i\xi}}_{-\pi}^\pi = \frac{e^{-2\pi i\xi\pi} - e^{2\pi i\xi\pi}}{-2\pi i\xi} = \frac{\sin(2\pi^2\xi)}{\pi\xi}.$$

\section*{Problem 12}
$$\widehat{f}(\xi) = \int_{-\infty}^\infty e^{-2\pi i\xi x}e^{-\pi x^2}dx = \int_{-\infty}^\infty e^{-2\pi i\xi x -\pi x^2}dx.$$
Completing the square:
$$-2\pi i\xi x -\pi x^2 = -\pi(x^2 + 2i\xi x),$$
$$x^2 + 2i\xi x = (x + i\xi)^2 - \xi^2,$$
$$-\pi(x^2 + 2i\xi x) = -\pi(x + i\xi)^2 + \pi\xi^2.$$
Therefore,
$$\widehat{f}(\xi) = e^{-\pi\xi^2}\int_{-\infty}^\infty e^{-\pi(x + i\xi)^2}dx = e^{-\pi\xi^2}\int_{-\infty}^\infty e^{-\pi y^2}dy$$
$$= e^{-\pi\xi^2}\int_{-\infty}^\infty e^{-(\sqrt\pi y)^2}dy = e^{-\pi\xi^2}\cdot \frac{1}{\sqrt{\pi}}\int_{-\infty}^\infty e^{-u^2}du = e^{-\pi\xi^2}.$$
For $f_t(x) = e^{-\pi tx^2} = f(\sqrt{t}x)$, we can use the scaling property:
$$g(x) = f(\lambda x) \implies \widehat{g}(\xi) = \frac{1}{\lambda}\widehat{f}(\xi/\lambda).$$
We get
$$\widehat{f_t}(\xi) = \frac{1}{\sqrt{t}}e^{-\pi\xi^2/t}.$$

\end{document}